% ============================================================
%  SCHOLA DOMESTICA — Magister Mathematica
%  Level 1 | Week 12 | Day 1
%  Topic: Subtraction Facts — −0 and −1 from 6 and Under
% ============================================================
\documentclass[12pt, letterpaper]{article}

\usepackage[top=1in, bottom=1in, left=1in, right=1in]{geometry}
\usepackage[T1]{fontenc}
\usepackage[utf8]{inputenc}
\usepackage{lmodern}
\usepackage{microtype}
\usepackage{amsmath}
\usepackage{xcolor}
\usepackage{tikz}
\usepackage{array}
\usepackage{enumitem}
\usepackage{multicol}
\usepackage{mdframed}
\usepackage{fancyhdr}

\definecolor{scholared}{RGB}{160,30,30}
\definecolor{scholagold}{RGB}{180,140,40}
\definecolor{scholacream}{RGB}{253,248,235}
\definecolor{scholagreen}{RGB}{50,110,60}
\definecolor{scholablue}{RGB}{40,80,140}

\pagestyle{fancy}
\fancyhf{}
\renewcommand{\headrulewidth}{0.6pt}
\fancyhead[L]{\small\textcolor{scholared}{\textbf{Schola Domestica — Magister Mathematica}}}
\fancyhead[R]{\small\textcolor{scholared}{Level 1 · Week 12 · Day 1}}
\fancyfoot[C]{\small\textcolor{gray}{— \thepage\ —}}

\newcommand{\answerline}{\underline{\hspace{2.5cm}}}
\newcommand{\biganswerline}{\underline{\hspace{4cm}}}
\newcommand{\numbox}{\tikz{\draw[scholablue,thick](0,0)rectangle(1.2cm,0.85cm);}}

% crossed-out dot
\newcommand{\xdot}{%
  \begin{tikzpicture}[baseline=-0.5ex]
    \filldraw[scholared!40] (0,0) circle (0.14cm);
    \draw[scholared, thick] (-0.12,-0.12) -- (0.12,0.12);
    \draw[scholared, thick] (-0.12,0.12) -- (0.12,-0.12);
  \end{tikzpicture}%
}
% plain dot
\newcommand{\odot}{\begin{tikzpicture}[baseline=-0.5ex]\filldraw[scholablue](0,0)circle(0.14cm);\end{tikzpicture}}

\newmdenv[
  backgroundcolor=scholacream,
  linecolor=scholared,
  linewidth=1.5pt,
  roundcorner=6pt,
  innertopmargin=8pt,
  innerbottommargin=8pt,
  innerleftmargin=10pt,
  innerrightmargin=10pt
]{lessonbox}

\newmdenv[
  backgroundcolor={rgb,255:red,235;green,245;blue,235},
  linecolor=scholagreen,
  linewidth=1.2pt,
  roundcorner=4pt,
  innertopmargin=6pt,
  innerbottommargin=6pt,
  innerleftmargin=8pt,
  innerrightmargin=8pt
]{enrichbox}

\begin{document}

\begin{center}
  {\LARGE\textbf{\textcolor{scholared}{Subtraction Facts:}}}\\[4pt]
  {\Large\textcolor{scholagold}{\textit{Taking Away Zero and One}}}\\[6pt]
  {\small Name: \biganswerline \hspace{1cm} Date: \answerline}
\end{center}

\vspace{4pt}
\noindent\textcolor{scholared}{\rule{\linewidth}{1pt}}
\vspace{4pt}

% IMAGE PROMPT (Beatrix Potter watercolour style):
% A little wren removes one shiny pebble from a row of five on a mossy log and tucks it under
% its wing. Soft green forest light, the remaining four pebbles glisten. No text, white background.

\begin{lessonbox}
\textbf{Subtraction} means \emph{taking away}. We use the $-$ sign.\\[6pt]
\begin{itemize}[leftmargin=1.5em, itemsep=4pt]
  \item \textbf{$-0$}: take away nothing — the number stays the same.\\
    \hspace{1em} $5 - 0 = 5$ \quad (nothing leaves)
  \item \textbf{$-1$}: count back one — one fewer.\\
    \hspace{1em} $5 - 1 = 4$ \quad (start at 5, go back 1: \textit{four})
\end{itemize}
\medskip
\noindent\textit{To show subtraction with dots: draw the whole group, then \textbf{cross out} the ones taken away.}
\end{lessonbox}

\bigskip

% ===== PART I: −0 FACTS =====
\noindent\textcolor{scholared}{\large\textbf{Part I — Taking Away Zero}} \hfill \textit{(Problems 1–4)}

\medskip
\noindent\textit{Cross out zero dots. Count what remains. Write the difference.}

\bigskip

\noindent
\begin{tabular}{@{}p{7.2cm} p{7.2cm}@{}}
% 1: 4−0=4
\textbf{1.}\quad
\begin{tikzpicture}[baseline=-0.5ex]
  \foreach \i in {1,...,4} { \filldraw[scholablue](\i*0.42,0)circle(0.14cm); }
\end{tikzpicture}
\quad $4 - 0 =$ \numbox
&
% 2: 6−0=6
\textbf{2.}\quad
\begin{tikzpicture}[baseline=-0.5ex]
  \foreach \i in {1,...,6} { \filldraw[scholablue](\i*0.38,0)circle(0.13cm); }
\end{tikzpicture}
\quad $6 - 0 =$ \numbox \\[16pt]
% 3: 3−0=3
\textbf{3.}\quad $3 - 0 =$ \numbox
&
% 4: 5−0=5
\textbf{4.}\quad $5 - 0 =$ \numbox
\end{tabular}

\bigskip\bigskip

% ===== PART II: −1 WITH CROSS-OUT =====
\noindent\textcolor{scholared}{\large\textbf{Part II — Taking Away One}} \hfill \textit{(Problems 5–10)}

\medskip
\noindent\textit{Cross out one dot (marked in red). Count what is left. Write the difference.}

\bigskip

\noindent
\begin{tabular}{@{}p{7.2cm} p{7.2cm}@{}}
% 5: 2−1=1
\textbf{5.}\quad
\begin{tikzpicture}[baseline=-0.5ex]
  \filldraw[scholablue](0.42,0)circle(0.14cm);
  % crossed-out dot
  \filldraw[scholared!40](0.9,0)circle(0.14cm);
  \draw[scholared,thick](0.77,-0.13)--(1.03,0.13);
  \draw[scholared,thick](0.77,0.13)--(1.03,-0.13);
\end{tikzpicture}
\quad $2 - 1 =$ \numbox
&
% 6: 3−1=2
\textbf{6.}\quad
\begin{tikzpicture}[baseline=-0.5ex]
  \foreach \i in {1,2} { \filldraw[scholablue](\i*0.42,0)circle(0.14cm); }
  \filldraw[scholared!40](1.32,0)circle(0.14cm);
  \draw[scholared,thick](1.19,-0.13)--(1.45,0.13);
  \draw[scholared,thick](1.19,0.13)--(1.45,-0.13);
\end{tikzpicture}
\quad $3 - 1 =$ \numbox \\[18pt]

% 7: 4−1=3
\textbf{7.}\quad
\begin{tikzpicture}[baseline=-0.5ex]
  \foreach \i in {1,2,3} { \filldraw[scholablue](\i*0.42,0)circle(0.14cm); }
  \filldraw[scholared!40](1.74,0)circle(0.14cm);
  \draw[scholared,thick](1.61,-0.13)--(1.87,0.13);
  \draw[scholared,thick](1.61,0.13)--(1.87,-0.13);
\end{tikzpicture}
\quad $4 - 1 =$ \numbox
&
% 8: 5−1=4
\textbf{8.}\quad
\begin{tikzpicture}[baseline=-0.5ex]
  \foreach \i in {1,...,4} { \filldraw[scholablue](\i*0.42,0)circle(0.14cm); }
  \filldraw[scholared!40](2.16,0)circle(0.14cm);
  \draw[scholared,thick](2.03,-0.13)--(2.29,0.13);
  \draw[scholared,thick](2.03,0.13)--(2.29,-0.13);
\end{tikzpicture}
\quad $5 - 1 =$ \numbox \\[18pt]

% 9: 6−1=5
\textbf{9.}\quad $6 - 1 =$ \numbox
&
% 10: 1−1=0
\textbf{10.}\quad $1 - 1 =$ \numbox \quad \textit{(what is left?)}
\end{tabular}

\bigskip\bigskip

% ===== PART III: NUMBER LINE SUBTRACTION =====
\noindent\textcolor{scholared}{\large\textbf{Part III — Count Back on the Number Line}} \hfill \textit{(Problems 11–14)}

\medskip
\noindent\textit{Start at the first number. Hop \textbf{left} (back). Write the difference.}

\bigskip

\noindent
\begin{tikzpicture}
  \draw[thick, scholablue] (0,0) -- (7.8,0);
  \foreach \n in {0,...,6} {
    \draw[thick] (\n*1.2+0.3,-0.12)--(\n*1.2+0.3,0.12);
    \node[below] at (\n*1.2+0.3,-0.18) {\small\textbf{\n}};
  }
\end{tikzpicture}

\bigskip

\noindent
\begin{tabular}{@{}p{3.5cm} p{3.5cm} p{3.5cm} p{3.5cm}@{}}
\textbf{11.}\; $5 - 1 =$ \numbox &
\textbf{12.}\; $4 - 0 =$ \numbox &
\textbf{13.}\; $6 - 1 =$ \numbox &
\textbf{14.}\; $3 - 1 =$ \numbox
\end{tabular}

\bigskip

\begin{enrichbox}
\textbf{\textcolor{scholagreen}{Addition and Subtraction are related!}}\\[3pt]
If $4 + 1 = 5$, then $5 - 1 = 4$.\\
If $3 + 0 = 3$, then $3 - 0 = 3$.\\[3pt]
\textit{Every addition fact has a subtraction partner. This week we will find them all!}
\end{enrichbox}

\vspace{0.3cm}
\noindent\textcolor{gray}{\small\textit{Ray's Primary Arithmetic — subtraction as the inverse of addition.}}

\end{document}
